\chapter*{Glosario}
\addcontentsline{toc}{chapter}{Glosario}
\markboth{GLOSARIO}{GLOSARIO}

\begin{description}
  \item \textbf{Framework:} Un framework es un conjunto de herramientas, bibliotecas y convenciones que proporcionan una estructura predefinida para el desarrollo de aplicaciones. Facilita la organización del código, la reutilización de componentes y la implementación de funcionalidades comunes, permitiendo a los desarrolladores centrarse en la lógica específica de su aplicación en lugar de preocuparse por detalles de bajo nivel.

  \item \textbf{JWT (JSON Web Token):} Es un estándar abierto (RFC 7519) que define un formato compacto y seguro para la transmisión de información entre partes como un objeto JSON. Se utiliza comúnmente para la autenticación y autorización en aplicaciones web, permitiendo a los servidores generar tokens que contienen información del usuario y sus permisos, los cuales pueden ser verificados y decodificados por el cliente para acceder a recursos protegidos.

  \item \textbf{Nodemailer:} Es una biblioteca de Node.js que facilita el envío de correos electrónicos desde aplicaciones Node.js. Proporciona una interfaz sencilla para configurar y enviar correos electrónicos utilizando diversos servicios de correo, como SMTP, Gmail, Outlook, entre otros

  \item \textbf{Supertest:} Es una biblioteca de pruebas para Node.js que permite realizar pruebas de integración en aplicaciones web. Proporciona una API fácil de usar para simular solicitudes HTTP y verificar las respuestas del servidor, lo que facilita la validación de rutas, controladores y middleware en aplicaciones Express u otros frameworks web.

  \item \textbf{Jest:} Es un framework de pruebas de JavaScript desarrollado por Facebook, diseñado para trabajar con proyectos que utilizan React, aunque también es compatible con otros entornos JavaScript. Jest proporciona una experiencia de prueba completa con características como pruebas unitarias, pruebas de integración, pruebas de extremo a extremo, y una interfaz de usuario amigable.

  \item \textbf{API:} Una API (Interfaz de Programación de Aplicaciones) es un conjunto de reglas y protocolos que permiten la comunicación entre diferentes software. Las APIs definen cómo los componentes de software deben interactuar, permitiendo a los desarrolladores acceder a funcionalidades específicas de una aplicación o servicio sin necesidad de conocer su implementación interna.

  \item \textbf{SMTP (Simple Mail Transfer Protocol):} Es un protocolo de comunicación utilizado para el envío de correos electrónicos a través de Internet. SMTP define las reglas y procedimientos para la transferencia de mensajes de correo electrónico entre servidores de correo, permitiendo a los usuarios enviar correos electrónicos desde sus clientes de correo electrónico a los servidores de correo y entre servidores de correo.

  \item \textbf{GitHub:} Es una plataforma de desarrollo colaborativo que utiliza el sistema de control de versiones Git. GitHub permite a los desarrolladores alojar, compartir y colaborar en proyectos de software, facilitando la gestión de código fuente, el seguimiento de cambios, la revisión de código y la integración continua.
\end{description}
